% ─────────────────────────────────────────────────────────────────────────────
\section{Dataset Description}
\label{sec:dataset}
% ─────────────────────────────────────────────────────────────────────────────

\subsection{NGAFID Overview}

The National General Aviation Flight Information Database (NGAFID) is a
real-world benchmark for aircraft health diagnosis research. It was collected
from a Cessna-172 fleet operated by an active flight school and correlates
23-dimensional sensor measurements with verified unscheduled maintenance records.
Table~\ref{tab:ngafid_overview} summarizes the key characteristics of the
dataset.

\begin{table}[H]
    \centering
    \caption{NGAFID dataset overview.}
    \label{tab:ngafid_overview}
    \begin{tabular}{ll}
        \toprule
        \textbf{Property} & \textbf{Value} \\
        \midrule
        Aircraft type         & Cessna-172 \\
        Operator              & Real flight school (anonymized) \\
        Total flights         & $>$28{,}000 (31{,}000 flight hours) \\
        Sensor channels       & 23 @ 1\,Hz \\
        Maintenance categories & 36 \\
        Unscheduled events    & $>$2{,}000 \\
        Preprocessing         & Forward-fill for NaN; cubic spline $\to$ 2{,}048 pts \\
        \midrule
        \multicolumn{2}{l}{\textit{Benchmark subset (2-day) — used in this work}} \\
        Total flights         & 11{,}446 \\
        Healthy flights       & 5{,}845 (51.1\%) \\
        Anomalous flights     & 5{,}601 (48.9\%) \\
        Fault classes (FC)    & 19 \\
        Cross-validation      & 5-fold stratified, physically isolated \\
        \bottomrule
    \end{tabular}
\end{table}

NGAFID provides two dataset configurations: the \textit{benchmark subset}
(19 classes, 11{,}446 flights) used for controlled evaluation, and the
\textit{overall dataset} (36 classes, 28{,}935 flights) for open-environment
assessment. This study uses the benchmark subset exclusively.

\subsection{Sensor Space}

The 23 sensor channels cover five critical aircraft subsystems.
Table~\ref{tab:sensors} lists all channels grouped by subsystem.

\begin{table}[H]
    \centering
    \caption{23-dimensional sensor space organized by subsystem.}
    \label{tab:sensors}
    \begin{tabular}{llc}
        \toprule
        \textbf{Subsystem} & \textbf{Channels} & \textbf{Count} \\
        \midrule
        Electrical   & \texttt{volt1}, \texttt{volt2}, \texttt{amp1}, \texttt{amp2} & 4 \\
        Fuel         & \texttt{FQtyL}, \texttt{FQtyR} & 2 \\
        Engine       & \texttt{E1 FFlow}, \texttt{E1 OilT}, \texttt{E1 OilP}, \texttt{E1 RPM} & 4 \\
        Cylinders (CHT) & \texttt{E1 CHT1}, \texttt{E1 CHT2}, \texttt{E1 CHT3}, \texttt{E1 CHT4} & 4 \\
        Cylinders (EGT) & \texttt{E1 EGT1}, \texttt{E1 EGT2}, \texttt{E1 EGT3}, \texttt{E1 EGT4} & 4 \\
        Flight state & \texttt{IAS}, \texttt{VSpd}, \texttt{AltMSL}, \texttt{NormAc} & 4 \\
        Environment  & \texttt{OAT} & 1 \\
        \midrule
        \textbf{Total} & & \textbf{23} \\
        \bottomrule
    \end{tabular}
\end{table}

All 23 channels are sampled at 1\,Hz. After preprocessing, each flight is
represented as a matrix $\mathbf{X} \in \mathbb{R}^{2048 \times 23}$, where
rows correspond to time steps and columns to sensor channels.

\subsection{Flight Duration Statistics}

Raw flight lengths vary considerably across the dataset, reflecting the diverse
operational modes of a flight school (student training, check rides, cross-country
flights). After interpolation to 2{,}048 fixed time steps, the original durations
are preserved as metadata:

\begin{itemize}
    \item \textbf{Minimum:} 605\,s ($\approx$\,10\,min)
    \item \textbf{Median:} 5{,}282\,s ($\approx$\,88\,min)
    \item \textbf{Maximum:} 20{,}679\,s ($\approx$\,345\,min)
\end{itemize}

\subsection{Class Distribution and Fault Taxonomy}

The 19 fault categories in the benchmark subset represent distinct maintenance
findings, each linked to a specific structural or functional component of the
aircraft. Table~\ref{tab:class_dist} lists all classes in descending order of
flight count.

\begin{table}[H]
    \centering
    \caption{Fault class distribution in the benchmark subset (5{,}601 anomalous flights).}
    \label{tab:class_dist}
    \begin{tabular}{clrr}
        \toprule
        \textbf{ID} & \textbf{Fault Category} & \textbf{Flights} & \textbf{\%} \\
        \midrule
        14 & Intake gasket leak / damage         & 2{,}097 & 37.4 \\
        18 & Rocker cover leak / loose / damage  & 1{,}024 & 18.3 \\
         2 & Baffle crack / damage / loose / miss &   304  &  5.4 \\
         4 & Intake tube / bolt / seal / boot loose or damage & 269 & 4.8 \\
         1 & Baffle plug need repair / replace    &   254  &  4.5 \\
         3 & Baffle screw missing / loose         &   211  &  3.8 \\
         6 & Baffle seal loose / damage           &   197  &  3.5 \\
         8 & Engine failure / fire / time out     &   161  &  2.9 \\
         5 & Baffle tie / tie rod loose or damage &   144  &  2.6 \\
        10 & Cylinder compression issue           &   143  &  2.6 \\
        11 & Engine run rough                     &   141  &  2.5 \\
         7 & Cylinder crack / fail / need repair  &   108  &  1.9 \\
        13 & Engine / propeller overspeed or damage &   94  &  1.7 \\
         9 & Engine idle / RPM issue              &    93  &  1.7 \\
        15 & Engine need repair / reinstall / clean &   80  &  1.4 \\
        12 & Engine seal / tube / bolt loose or damage &  76 & 1.4 \\
        16 & Oil cooler need maintenance          &    75  &  1.3 \\
        17 & Pilot / in-flight noticed issue      &    75  &  1.3 \\
         0 & Baffle rivet loose / missing / damage &   55  &  1.0 \\
        \midrule
        \textbf{---} & \textbf{Total anomalous} & \textbf{5{,}601} & \textbf{100.0} \\
        \bottomrule
    \end{tabular}
\end{table}

The imbalance ratio between the largest class (intake gasket, $n$\,=\,2{,}097)
and the smallest class (baffle rivet, $n$\,=\,55) is \textbf{38.1:1}, creating
an extreme long-tailed distribution that poses a significant challenge for
minority class recognition.

\subsection{Key Data Challenges}
\label{subsec:challenges}

NGAFID exhibits three fundamental characteristics that distinguish it from
synthetic benchmarks and directly constrain the achievable performance of any
classification approach \cite{convtok2023, lmsd2025}.

\subsubsection{Target Component Sparsity}

When the feature matrix $\mathbf{X}$ is projected into principal component space,
fault-discriminative information concentrates in low-variance components (notably
PC9 and PC19--23 by per-component $F$-statistic), which each account for
\textbf{less than 0.1\% of total variance}. The dominant high-variance components
(PC1--PC2) capture operational variability common to both healthy and anomalous
flights---altitude profile, throttle inputs, ambient temperature effects---and are
nearly non-discriminative, masking the fault signatures that classifiers need to
detect.

Empirically, the Repr.~B feature matrix (184 features) requires \textbf{51
principal components} to explain 95\% of its variance. This high intrinsic
dimensionality relative to the number of fault-informative components explains
why deep models with local receptive fields outperform global statistical methods
on FC.

\subsubsection{Multi-Stage Coupling}

Each flight traverses multiple operationally heterogeneous phases: ground taxi,
takeoff and climb, cruise, descent and approach, and ground rollout. These phases
have statistically distinct dynamics for virtually every sensor channel---RPM,
fuel flow, CHT, and altitude all follow phase-specific profiles. This
heterogeneity creates strong temporal dependencies that prohibit phase-wise
isolation; a feature vector that captures only one flight phase misses the
fault's temporal context. Conversely, a global feature vector aggregates across
all phases, diluting phase-specific anomalies.

\subsubsection{Dual Imbalance}

NGAFID exhibits imbalance at two levels simultaneously:
\begin{itemize}
    \item \textbf{Quantitative (AD level):} The benchmark subset is nearly
          balanced between healthy (51.1\%) and anomalous (48.9\%) flights.
          However, the \textit{overall} NGAFID dataset has a substantial
          majority of healthy flights, reflecting real fleet operations where
          chronic wear dominates over acute failures.
    \item \textbf{Structural (FC level):} Among the 19 fault classes, the
          distribution is severely long-tailed with an imbalance ratio of 38.1:1.
          The two head classes alone (intake gasket + rocker cover) account for
          55.7\% of all anomalous flights. This reflects \textit{objective fleet
          reality}---some faults are inherently more prevalent---rather than
          sampling bias.
\end{itemize}

\subsection{Cross-Validation Protocol}
\label{subsec:cv}

All experiments use a \textbf{5-fold stratified cross-validation} scheme that
maintains class proportions within each fold while enforcing strict physical
isolation: no two flights from the same physical aircraft appear in both the
training and test partitions of the same fold. This prevents optimistic
performance estimates that would arise from data leakage between temporally
correlated flights.

Fold composition for the AD task is shown in Table~\ref{tab:cv_folds}.

\begin{table}[H]
    \centering
    \caption{5-fold CV composition for the AD task.}
    \label{tab:cv_folds}
    \begin{tabular}{ccc}
        \toprule
        \textbf{Fold} & \textbf{Healthy} & \textbf{Anomalous} \\
        \midrule
        0 & 1{,}167 & 1{,}123 \\
        1 & 1{,}123 & 1{,}166 \\
        2 & 1{,}162 & 1{,}127 \\
        3 & 1{,}192 & 1{,}097 \\
        4 & 1{,}201 & 1{,}088 \\
        \midrule
        \textbf{Total} & \textbf{5{,}845} & \textbf{5{,}601} \\
        \bottomrule
    \end{tabular}
\end{table}

Feature normalization (z-score standardization) is applied within each fold,
with mean and standard deviation estimated exclusively from the training split
and applied to the test split to prevent test-set contamination.
